\answer{ex:02:1}{阵列长 $3.2$~mm，间距 $25\um$：$\approx129$ 个检测器。发放的是宽 $v\,\tau\ln2=1.7$~mm 的一条带，$\approx69$ 个检测器，超过阵列的一半；方向取这条带的中心。}
\answer{ex:02:2}{窗口为 $-\tau\ln\frac{w_2}{\theta-w_1}\le\delta\le\tau\ln\frac{w_1}{\theta-w_2}$，即 $[-0.235,\,0.178]\ms$；中心 $c=\frac\tau2\ln\frac{w_1(\theta-w_1)}{w_2(\theta-w_2)}=-28\,\mu$s。方向向较强的耳朵一侧偏移 $28\,\mu$s，接近三个分辨步长。}
\answer{ex:02:3}{需要满足 $s_iw_{ij}s_j\ge0$ 的 $s_i$；当且仅当所有环都是偶环时它们存在。相互抑制可以化为这种网络（没有持续振荡），\nt{GLUT}--\nt{GABA} 环则不能（可能振荡）。}
\answer{ex:02:4}{六个神经元 $g_k=[x+y+z\ge k]$ 和 $\bar g_k=[\bar x+\bar y+\bar z\ge4-k]$，$k=1,2,3$；$C=g_2$，$\bar C=\bar g_2$，$S=[g_1+\bar g_2+g_3\ge2]$，$\bar S=[\bar g_1+g_2+\bar g_3\ge2]$：共八个神经元。用与、或元件需 $12$ 个。}
\answer{ex:02:5}{$T_{\min}(0.6)=0.718\,\tau$；$I_c=0.225$，在阈值附近 $T_{\min}\propto(I-I_c)^{-1/2}$。}
\answer{ex:02:6}{(a)~$500$ 环，$5\cdot10^4$ 个神经元；散布 $4.5\ms$（$0.45\,\%$），相对误差 $\propto T^{-1/2}$。(b)~所有环共有的随机延迟因子，散布为 $5\,\%$。}
\answer{ex:02:7}{每 $\tau_F\ln2=0.69\,\text{s}$ 刷新一次：每个神经元每秒 $4.3$ 个脉冲，而触发器是 $20$ 个，节省 $4.6$ 倍（每个脉冲 $10^{-10}$~J 时，$4\cdot10^{-7}$~J 对 $2\cdot10^{-6}$~J）。触发器丢失比特；电容器只要压制开始于刷新后不超过 $0.49\,\text{s}$，就能保住比特（概率 $0.71$）。}
